\documentclass[12pt,a4paper]{article}
\usepackage[UTF8]{ctex}
\usepackage{geometry}
\usepackage{graphicx}
\usepackage{booktabs}
\usepackage{amsmath}
\usepackage{amssymb}
\usepackage{float}
\usepackage{caption}
\usepackage{subcaption}
\usepackage{hyperref}
\usepackage{xcolor}
\usepackage{fancyhdr}

% 页面设置
\geometry{left=2.5cm,right=2.5cm,top=2.5cm,bottom=2.5cm}

% 页眉页脚
\pagestyle{fancy}
\fancyhf{}
\fancyhead[L]{BCI实验报告}
\fancyhead[R]{\today}
\fancyfoot[C]{\thepage}

% 标题格式
\usepackage{titlesec}
\titleformat{\section}{\Large\bfseries\color{blue!60!black}}{}{0em}{}
\titleformat{\subsection}{\large\bfseries\color{blue!40!black}}{}{0em}{}

% 代码环境
\usepackage{listings}
\lstset{
    basicstyle=\ttfamily\small,
    breaklines=true,
    frame=single,
    backgroundcolor=\color{gray!10},
    language=Python
}

\begin{document}

% 标题页
\begin{titlepage}
    \centering
    \vspace*{2cm}
    
    {\Huge\bfseries {{ experiment_name }}\par}
    \vspace{1cm}
    {\Large BCI运动想象实验报告\par}
    \vspace{2cm}
    
    \begin{tabular}{rl}
        实验ID: & {{ experiment_id }} \\
        研究者: & {{ researcher }} \\
        机构: & {{ institution }} \\
        实验日期: & {{ experiment_date }} \\
        报告生成时间: & {{ generation_time }} \\
    \end{tabular}
    
    \vfill
    {\large BCI运动想象康复训练平台}
\end{titlepage}

% 目录
\tableofcontents
\newpage

% 正文开始
\section{实验基本信息}

\subsection{实验概述}
\textbf{实验描述}: {{ description }}

\subsection{被试信息}
\begin{table}[H]
    \centering
    \begin{tabular}{ll}
        \toprule
        项目 & 内容 \\
        \midrule
        被试编号 & {{ subject_id }} \\
        年龄 & {{ age }}岁 \\
        性别 & {{ gender }} \\
        利手 & {{ handedness }} \\
        经验水平 & {{ experience_level }} \\
        纳入标准 & {{ inclusion_criteria }} \\
        \bottomrule
    \end{tabular}
    \caption{被试信息}
\end{table}

\section{数据采集与预处理}

\subsection{数据采集参数}
\begin{table}[H]
    \centering
    \begin{tabular}{ll}
        \toprule
        参数 & 值 \\
        \midrule
        设备型号 & {{ device_name }} \\
        通道数 & {{ channel_count }} \\
        采样率 & {{ sampling_rate }} Hz \\
        参考电极 & {{ reference }} \\
        接地电极 & {{ ground }} \\
        阻抗阈值 & {{ impedance_threshold }} k$\Omega$ \\
        \bottomrule
    \end{tabular}
    \caption{数据采集参数}
\end{table}

\subsection{预处理参数}
\begin{table}[H]
    \centering
    \begin{tabular}{ll}
        \toprule
        参数 & 值 \\
        \midrule
        带通滤波 & {{ bandpass_low }}-{{ bandpass_high }} Hz \\
        陷波滤波 & {{ notch_frequency }} Hz \\
        伪迹去除方法 & {{ artifact_removal_method }} \\
        眼电通道 & {{ eog_channels }} \\
        坏道剔除 & {{ bad_channels }} \\
        \bottomrule
    \end{tabular}
    \caption{预处理参数}
\end{table}

\section{特征提取与分类}

\subsection{特征提取参数}
\begin{table}[H]
    \centering
    \begin{tabular}{ll}
        \toprule
        参数 & 值 \\
        \midrule
        提取方法 & {{ feature_method }} \\
        时间窗口 & {{ time_window[0] }}-{{ time_window[1] }}秒 \\
        频带 & {{ frequency_bands }} \\
        成分数量 & {{ n_components }} \\
        \bottomrule
    \end{tabular}
    \caption{特征提取参数}
\end{table}

\subsection{分类结果}

\begin{table}[H]
    \centering
    \begin{tabular}{lcc}
        \toprule
        指标 & 值 & 标准差 \\
        \midrule
        准确率 & {{ "%.1f"|format(accuracy*100) }}\% & $\pm${{ "%.1f"|format(accuracy_std*100) }}\% \\
        F1-Score & {{ "%.1f"|format(f1_score*100) }}\% & $\pm${{ "%.1f"|format(f1_std*100) }}\% \\
        Kappa系数 & {{ "%.3f"|format(kappa) }} & $\pm${{ "%.3f"|format(kappa_std) }} \\
        AUC & {{ "%.3f"|format(auc) }} & $\pm${{ "%.3f"|format(auc_std) }} \\
        \bottomrule
    \end{tabular}
    \caption{分类性能指标}
\end{table}

\subsubsection{混淆矩阵}
\begin{table}[H]
    \centering
    \begin{tabular}{lcc}
        \toprule
         & 预测左手 & 预测右手 \\
        \midrule
        实际左手 & {{ cm[0,0] }} & {{ cm[0,1] }} \\
        实际右手 & {{ cm[1,0] }} & {{ cm[1,1] }} \\
        \bottomrule
    \end{tabular}
    \caption{混淆矩阵}
\end{table}

\subsubsection{算法配置}
\begin{table}[H]
    \centering
    \begin{tabular}{ll}
        \toprule
        项目 & 值 \\
        \midrule
        算法 & {{ algorithm }} \\
        交叉验证 & {{ cv_method }} ({{ n_folds }}折) \\
        训练时间 & {{ training_time }}秒 \\
        推理时间 & {{ inference_time }}毫秒 \\
        \bottomrule
    \end{tabular}
    \caption{算法配置信息}
\end{table}

\section{统计分析}

{% for test in statistical_tests %}
\subsection{{ test.test_name }}
\begin{table}[H]
    \centering
    \begin{tabular}{ll}
        \toprule
        统计量 & {{ "%.4f"|format(test.statistic) }} \\
        p值 & {{ "%.4f"|format(test.p_value) }} \\
        显著性 & {% if test.significant %}显著 ($p < 0.05$){% else %}不显著 ($p \ge 0.05$){% endif %} \\
        效应量 & {{ "%.4f"|format(test.effect_size) }} \\
        95\%置信区间 & ({{ "%.4f"|format(test.confidence_interval[0]) }}, {{ "%.4f"|format(test.confidence_interval[1]) }}) \\
        比较组 & {{ test.comparison_groups[0] }} vs {{ test.comparison_groups[1] }} \\
        \bottomrule
    \end{tabular}
    \caption{{ test.test_name }}结果
\end{table}
{% endfor %}

\section{可视化结果}

{% for figure_name, figure_path in figures.items() %}
\subsection{{ figure_name }}
\begin{figure}[H]
    \centering
    \includegraphics[width=0.8\textwidth]{{ figure_path }}
    \caption{{ figure_name }}
    \label{fig:{{ figure_name|replace(" ", "_") }}}
\end{figure}
{% endfor %}

\section{结论与讨论}

\subsection{主要发现}
% 自动生成的主要发现
{% if accuracy > 0.85 %}
实验结果表明，{{ algorithm }}算法在运动想象分类任务上取得了良好的性能（准确率 = {{ "%.1f"|format(accuracy*100) }}\%），显著优于随机水平。
{% else %}
实验结果显示分类性能有待进一步提升，建议优化特征提取方法或增加训练数据。
{% endif %}

\subsection{局限性}
% 自动识别的局限性
{% if n_subjects < 20 %}
\textbf{样本量限制}：当前实验仅包含{{ n_subjects }}名被试，样本量较小，可能影响结果的普适性。
{% endif %}

{% if cv_folds < 10 %}
\textbf{交叉验证}：采用{{ cv_folds }}折交叉验证，建议使用更多折数以获得更稳定的估计。
{% endif %}

\subsection{未来工作}
\begin{itemize}
    \item 增加被试数量，提高统计功效
    \item 探索更先进的深度学习模型
    \item 实现实时反馈闭环系统
    \item 开展多中心临床验证研究
\end{itemize}

\section*{参考文献}
\begin{enumerate}
    \item Zhang, X., et al. (2024). EEGNet: A Compact Convolutional Network for EEG-based Brain-Computer Interfaces. \textit{IEEE TBME}.
    \item Lawhern, V. J., et al. (2018). EEGNet: A compact convolutional neural network for EEG-based brain-computer interfaces. \textit{Journal of Neural Engineering}.
    \item Schirrmeister, R. T., et al. (2017). Deep learning with convolutional neural networks for EEG decoding and visualization. \textit{Human Brain Mapping}.
\end{enumerate}

\section*{附录}
\subsection{原始数据说明}
所有原始数据已按照数据管理规范存储，可通过以下方式访问：
\begin{itemize}
    \item 数据存储路径: \texttt{./datasets/{{ experiment_id }}}
    \item 预处理代码: \texttt{./src/preprocessing/}
    \item 分析脚本: \texttt{./src/pipeline/}
\end{itemize}

\subsection{软件环境}
\begin{lstlisting}
操作系统: {{ os_info }}
Python版本: {{ python_version }}
主要依赖:
- numpy: {{ numpy_version }}
- scipy: {{ scipy_version }}
- scikit-learn: {{ sklearn_version }}
- torch: {{ torch_version }}
- mne: {{ mne_version }}
\end{lstlisting}

\newpage
\centerline{\small 本报告由BCI运动想象康复训练平台自动生成}
\centerline{\small 报告生成时间: {{ generation_time }}}

\end{document}